\section{A.L. 1:Corriente de saturación $I_{DSS}$}

	\sangria{}En esta actividad buscamos conocer experimentalmente el valor de $I_{DSS}$, el cual lo podemos obtener poniedo la terminal gate a tierra para que la tensión $V_{GS}$ sea 0 y la corriente que circule por el canal n sea la máxima.

	\begin{tikzpicture}
	% Paths, nodes and wires:
	\node[njfet, tr circle](N1) at (4, 4.25){} node[anchor=west] at (N1.text){$J1 - MPF102$};
	\draw (4, 5) -- (7, 5) -| (7, 4);
	\draw (7, 4) to[american voltage source, l={$V$}] (7, 2.25);
	\node[ground] at (5, 2){};
	\draw (4, 3.48) -| (5, 2);
	\draw (7, 2.25) |- (5, 2);
	\draw (3.02, 3.98) |- (5, 2);
\end{tikzpicture}

	\subsection{Simulación}
	\imagen[Simulación Actividad 1]{6cm}{./imagenes/Simulacion_1.png}
	\imagen[Simulación Grafica Actividad 1]{6cm}{./imagenes/Simulacion_1_grafica.png}

	\sangria{}En la simulación podemos ver que $V_p \approx 12.3V$, en la simulación incluimos una resistencia $R_D$ para no quemar el transistor, el cálculo
